% ==========================================================================
%  Chapter 75 -- Concurrent Multiscale Architecture:
%                From FE² to a Publishable Framework in fall_n
% ==========================================================================

\chapter{Concurrent Multiscale Architecture: Diagnostic, Design, and Roadmap}
\label{ch:multiscale-architecture}


% =========================================================================
\section{Introduction}
\label{sec:ms-arch-intro}

The previous chapters developed the building blocks of a concurrent
multiscale framework for reinforced-concrete structures: a prismatic
domain builder (Chapter~\ref{ch:prismatic-domain-ho}), a persistent
nonlinear sub-model solver (Chapter~\ref{ch:nl-sub-model-evolver}),
arc-length adaptive control (Chapter~\ref{ch:arc-length}), and
two-way coupling via homogenised tangent injection
(Chapter~\ref{ch:two-way-fe2}).  This chapter steps back to perform a
comprehensive architectural audit of the current implementation,
identifies rigidities and design debts, and lays out a detailed
roadmap toward a flexible, SOLID-conformant, HPC-ready multiscale
analysis framework suitable for publication and external use.

The discussion is organised as follows.
Section~\ref{sec:ms-fe2-review} reviews the classical
FE\textsuperscript{2} method of Feyel and Chaboche and explains how
our implementation adapts their ideas.
Section~\ref{sec:ms-current-status} audits the current code against
SOLID principles.
Section~\ref{sec:ms-proposed-arch} presents the target architecture.
Section~\ref{sec:ms-reuse} catalogues reusable and hard-coded code.
Section~\ref{sec:ms-hpc} addresses parallelism and HPC scalability.
Section~\ref{sec:ms-plan} provides a phased implementation plan.


% =========================================================================
\section{The Classical FE\textsuperscript{2} Method}
\label{sec:ms-fe2-review}

% -------------------------------------------------------------------------
\subsection{Origin and Formulation}

The FE\textsuperscript{2} method was formalised by
Feyel \cite{Feyel1999} and subsequently refined by
Feyel and Chaboche \cite{FeyelandChaboche2000} as a systematic
framework for concurrent multiscale analysis.  The fundamental idea
is to replace the phenomenological constitutive law at each Gauss
point of a macro-scale finite-element model by a boundary-value
problem solved on a Representative Volume Element (RVE) at the
micro-scale.  The method proceeds as follows:

\begin{enumerate}[label=(\roman*)]
  \item \textbf{Localisation.}  At a given macro-scale Gauss point
    with deformation gradient $\bar{\bm{F}}$ (or strain
    $\bar{\bm{\varepsilon}}$ in the small-strain case), the macro
    kinematic field is imposed as boundary conditions on the RVE
    domain $\Omega_\mu$:
    \begin{equation}
      \label{eq:fe2-bc}
      \mathbf{u}(\mathbf{x}) =
        \bar{\bm{\varepsilon}} \cdot \mathbf{x}
        + \tilde{\mathbf{u}}(\mathbf{x}),
      \qquad \mathbf{x} \in \partial\Omega_\mu,
    \end{equation}
    where $\tilde{\mathbf{u}}$ is a fluctuation field satisfying
    appropriate boundary conditions (linear displacement, periodic,
    or minimal kinematic).

  \item \textbf{Micro-scale solve.}  The RVE boundary-value problem
    is solved with the micro-scale material laws (plasticity, damage,
    cracking, etc.) to obtain the local stress field
    $\bm{\sigma}(\mathbf{x})$.

  \item \textbf{Homogenisation.}  The macro-scale stress is obtained
    by volume averaging over the RVE:
    \begin{equation}
      \label{eq:fe2-homog}
      \bar{\bm{\sigma}} = \frac{1}{|\Omega_\mu|}
        \int_{\Omega_\mu} \bm{\sigma}(\mathbf{x}) \; \mathrm{d}\Omega.
    \end{equation}

  \item \textbf{Consistent tangent.}  The macro-scale algorithmic
    tangent $\bar{\mathbb{C}} = \partial\bar{\bm{\sigma}}
    / \partial\bar{\bm{\varepsilon}}$ is computed either by static
    condensation of the micro-scale tangent stiffness or by numerical
    perturbation.  Feyel \cite{Feyel1999} derived the condensed
    tangent:
    \begin{equation}
      \label{eq:fe2-cond-tangent}
      \bar{\mathbb{C}} = \frac{1}{|\Omega_\mu|}
        \left(
          \mathbf{K}_{bb} -
          \mathbf{K}_{bi} \, \mathbf{K}_{ii}^{-1} \, \mathbf{K}_{ib}
        \right),
    \end{equation}
    where the subscripts $b$ and $i$ denote boundary and interior
    degrees of freedom, respectively.
\end{enumerate}

\noindent
The key contribution of Feyel and Chaboche was demonstrating that
this concurrent scheme --- nested FE solves --- could be implemented
within a single FE code, preserving the algorithmic tangent's
consistency and thereby maintaining quadratic convergence of the
macro-scale Newton iteration.  They applied the method to
metal--matrix composites with elasto-viscoplastic constituents,
showing that the macro-scale response automatically captures complex
nonlinear phenomena (damage localisation, strain gradients) that
phenomenological models struggle to represent.

% -------------------------------------------------------------------------
\subsection{Extensions in the Literature}

Since the original formulation, FE\textsuperscript{2} has been
extended in numerous directions:

\begin{itemize}
  \item \textbf{Periodic boundary conditions} for textile composites
    (Miehe et al.\ \cite{Miehe2002}).
  \item \textbf{Computational homogenisation for shells and beams}
    using dimensional reduction at the micro-scale
    (Geers et al.\ \cite{Geers2010}).
  \item \textbf{Second-order homogenisation} incorporating macro-scale
    strain gradients (Kouznetsova et al.\ \cite{Kouznetsova2004}).
  \item \textbf{Adaptive FE\textsuperscript{2}} where the micro-scale
    solve is activated only at Gauss points exceeding a damage
    threshold (Ghosh et al.\ \cite{Ghosh2001}).
  \item \textbf{HPC implementations} exploiting the embarrassingly
    parallel nature of the micro-scale solves (Feyel \cite{Feyel2003},
    Schröder \cite{Schroeder2014}).
\end{itemize}

% -------------------------------------------------------------------------
\subsection{Adaptation in \texttt{fall\_n}}
\label{sec:ms-fe2-adaptation}

Our implementation in \texttt{fall\_n} adapts the original
FE\textsuperscript{2} framework with several domain-specific
modifications tailored to structural engineering of RC frames under
seismic loading:

\begin{enumerate}[label=(\arabic*)]
  \item \textbf{Beam--continuum scale bridge.}
    Unlike the classical FE\textsuperscript{2}, where both scales
    operate in the same dimensionality (3D--3D or 2D--2D), our
    macro-scale is a one-dimensional beam element (Timoshenko 3D with
    fibre sections) and the micro-scale is a three-dimensional
    continuum RVE.  The localisation step
    (\ref{eq:fe2-bc}) is replaced by a \emph{Timoshenko kinematic
    field} that maps section deformations
    $\mathbf{e} = \{\varepsilon_0, \kappa_y, \kappa_z, \gamma_y,
    \gamma_z, \theta'\}$ to displacement boundary conditions on the
    prismatic RVE faces (see
    Chapter~\ref{ch:structural-scale-reconstruction}).

  \item \textbf{Boundary-reaction homogenisation.}
    Instead of volume-averaging the stress field (which requires
    access to the stored material state at all Gauss points), we
    extract section forces directly from the assembled
    internal-force vector at the RVE boundary face.  This approach
    (Section~\ref{sec:fe2-reactions}) is more robust for the
    perturbation tangent loop, as it evaluates the constitutive
    response from the displacement field without relying on stored
    material state that may be stale between non-committed SNES
    solves.

  \item \textbf{Damage-triggered transition.}
    The micro-scale solve is activated only at beam elements exceeding
    a damage threshold (first steel fiber yield), implementing the
    adaptive FE\textsuperscript{2} philosophy of Ghosh et al.\
    \cite{Ghosh2001}.  A \texttt{TransitionDirector} monitors damage
    at every time step and triggers the transition seamlessly.

  \item \textbf{Persistent material state evolution.}
    Unlike one-shot homogenisation where the RVE is solved from
    scratch at each step, our \texttt{NonlinearSubModelEvolver}
    maintains the full material state (crack history, plastic strains,
    damage variables) across all time steps.  This is essential for
    cyclic seismic loading, where unloading--reloading paths govern
    the hysteretic energy dissipation.

  \item \textbf{Mixed-model RVE.}
    The RVE employs a \texttt{MixedModel} combining
    \texttt{ContinuumElement} (KoBatheConcrete3D for the concrete
    matrix) and \texttt{TrussElement} (MenegottoPinto steel for
    embedded rebar), enabling a physically motivated representation
    of the reinforced cross-section.

  \item \textbf{Forward-difference perturbation tangent.}
    The consistent tangent (\ref{eq:fe2-cond-tangent}) requires
    access to the condensed micro-scale stiffness matrix, which is
    not readily available when using PETSc's SNES (Newton) solver as
    a black box.  Instead, we compute the $6 \times 6$ homogenised
    tangent $\mathbf{D}_{\text{hom}}$ by forward-difference
    perturbation of the section deformation components, using the
    boundary-reaction extraction at each perturbation step.
\end{enumerate}

\noindent
The principal contribution of our approach, relative to the classical
FE\textsuperscript{2}, is the \emph{beam--continuum dimensional
bridge} with \emph{boundary-reaction homogenisation}: by reducing
the upscaling from a $6n_{\text{dof}}$-dimensional condensation to a
six-component reaction summation, we achieve a computationally
efficient coupling that avoids the full stiffness condensation and
operates naturally within the PETSc solver framework.


% =========================================================================
\section{Current Implementation Status}
\label{sec:ms-current-status}

% -------------------------------------------------------------------------
\subsection{Component Inventory}

Table~\ref{tab:ms-components} summarises the current multiscale
components, their locations, and coupling directions.

\begin{table}[htbp]
  \centering
  \caption{Multiscale component inventory.}
  \label{tab:ms-components}
  \small
  \begin{tabular}{@{}lllc@{}}
    \toprule
    \textbf{Component} & \textbf{File} & \textbf{Role} & \textbf{Direction} \\
    \midrule
    \texttt{MultiscaleCoordinator}
      & \texttt{MultiscaleCoordinator.hh}
      & Kinematic downscaling orchestrator
      & $\downarrow$ \\
    \texttt{NonlinearSubModelEvolver}
      & \texttt{NonlinearSubModelEvolver.hh}
      & Persistent SNES-based RVE solver
      & $\updownarrow$ \\
    \texttt{section\_forces\_from\_reactions()}
      & idem
      & Boundary-reaction homogenisation
      & $\uparrow$ \\
    \texttt{compute\_homogenized\_tangent()}
      & idem
      & $6 \times 6$ perturbation tangent
      & $\uparrow$ \\
    \texttt{compute\_homogenized\_forces()}
      & idem
      & Section forces via \texttt{HomogenizedSection}
      & $\uparrow$ \\
    \texttt{set\_homogenized\_tangent/forces}
      & \texttt{BeamElement.hh}
      & Injection into beam section
      & $\uparrow$ \\
    Staggered iteration loop
      & \texttt{main\_lshaped\_multiscale.cpp}
      & Coupling iteration with relaxation
      & $\updownarrow$ \\
    \texttt{SubModelSolver}
      & \texttt{SubModelSolver.hh}
      & Stateless one-shot continuum solver
      & $\downarrow$ \\
    \texttt{HomogenizedSection}
      & \texttt{HomogenizedSection.hh}
      & Stress-to-resultants projection
      & $\uparrow$ \\
    \texttt{FieldTransfer / SubModelSpec}
      & \texttt{FieldTransfer.hh}
      & Kinematic field transfer
      & $\downarrow$ \\
    \bottomrule
  \end{tabular}
\end{table}

% -------------------------------------------------------------------------
\subsection{Verification Test Results}

The complete test suite (63 tests) passes after the
boundary-reaction tangent implementation.  The critical
\texttt{evolver\_advanced} test (21 sub-tests) reports:
\begin{align*}
  \mathbf{D}_{\text{hom}} \text{ diagonal}
    &= [1125.89, \; 6.51, \; 6.51, \; 72.09, \; 72.09, \; 3.52], \\
  \mathbf{f}_{\text{hom}}
    &= [0.114, \; 0, \; 0, \; {\sim}0, \; {\sim}0, \; 0],
\end{align*}
confirming all six diagonal entries of the tangent are positive and
the homogenised forces exhibit the expected axial-only response
under pure tension.

% -------------------------------------------------------------------------
\subsection{SOLID Principles Audit}
\label{sec:ms-solid-audit}

Despite passing all tests, the current implementation exhibits several
design debts when evaluated against the SOLID principles of
object-oriented design.

\paragraph{Single Responsibility Principle (SRP).}
\texttt{NonlinearSubModelEvolver} consolidates at least five distinct
responsibilities in a single class:
\begin{itemize}
  \item SNES-based nonlinear solving (residual/Jacobian callbacks,
    convergence control);
  \item Boundary condition transfer (writing imposed values from
    section kinematics);
  \item Material model construction (hard-coded
    \texttt{KoBatheConcrete3D} and \texttt{MenegottoPintoSteel});
  \item VTK post-processing (mesh, Gauss cloud, crack plane export);
  \item Crack tracking and summary statistics.
\end{itemize}
Each of these is an independently evolving concern.

\paragraph{Open/Closed Principle (OCP).}
The material models used in the RVE are hard-coded:
\begin{lstlisting}[caption={Hard-coded material in \texttt{first\_solve()}.}]
InelasticMaterial<KoBatheConcrete3D> mat_inst{fc_};
Material<Policy> mat{mat_inst, InelasticUpdate{}};
// ...
MenegottoPintoSteel steel_impl{rebar_E_, rebar_fy_, rebar_b_};
\end{lstlisting}
Extending the framework to alternative concrete models (e.g.\ CDPM2,
Mazars damage) or steel models (Ramberg--Osgood, explicit cyclic)
requires modifying the class internals.

\paragraph{Liskov Substitution Principle (LSP).}
\texttt{SubModelSolver} (stateless, per-step) and
\texttt{NonlinearSubModelEvolver} (stateful, persistent) do not
share an interface, despite providing overlapping functionality:
both solve a sub-model and extract results.  Client code cannot
substitute one for the other.

\paragraph{Interface Segregation Principle (ISP).}
The evolver exposes a monolithic interface that forces clients to
depend on crack tracking, VTK export, and arc-length control even
when they only require the solver or homogenisation functionality.

\paragraph{Dependency Inversion Principle (DIP).}
The staggered coupling loop in
\texttt{main\_lshaped\_multiscale.cpp} (approximately 80 lines)
depends directly on concrete types (\texttt{NonlinearSubModelEvolver},
\texttt{BeamElemT}, \texttt{COL\_B}, \texttt{COL\_H}).  The coupling
algorithm, convergence criterion, and relaxation strategy are all
embedded in the application driver rather than abstracted behind
stable interfaces.


% =========================================================================
\section{Proposed Architecture: \texttt{MultiscaleAnalysis}}
\label{sec:ms-proposed-arch}

% -------------------------------------------------------------------------
\subsection{Design Overview}

The target architecture decomposes the current monolith into
a layered system of abstractions, as shown in
Figure~\ref{fig:ms-arch}.

\begin{figure}[htbp]
  \centering
  \begin{tikzpicture}[
      box/.style={draw, rounded corners, minimum width=4.8cm,
                  minimum height=1cm, text centered,
                  font=\small\ttfamily},
      iface/.style={draw, dashed, rounded corners, minimum width=4cm,
                    minimum height=0.8cm, text centered,
                    font=\small\itshape},
      arr/.style={-{Stealth[length=3mm]}, thick},
      >=Stealth
    ]
    % Top level
    \node[box, fill=blue!10] (msa) at (0,0)
      {MultiscaleAnalysis};

    % Policies
    \node[iface, fill=yellow!10] (sbp) at (-4.5,-1.8)
      {ScaleBridgePolicy};
    \node[iface, fill=yellow!10] (ccp) at (0,-1.8)
      {CouplingConvergence};
    \node[iface, fill=yellow!10] (rlx) at (4.5,-1.8)
      {RelaxationPolicy};

    % MultiscaleModel
    \node[box, fill=green!10] (msm) at (0,-3.6)
      {MultiscaleModel};

    % Sub-components
    \node[box, fill=gray!10] (gm) at (-4,-5.2)
      {GlobalModel (beams)};
    \node[box, fill=gray!10] (lm) at (4,-5.2)
      {LocalModelAdapter[]};

    % Bottom
    \node[iface, fill=orange!10] (mf) at (-4,-7)
      {MaterialFactory};
    \node[iface, fill=orange!10] (ht) at (4,-7)
      {HomogenisationStrategy};

    % Arrows
    \draw[arr] (msa) -- (sbp);
    \draw[arr] (msa) -- (ccp);
    \draw[arr] (msa) -- (rlx);
    \draw[arr] (msa) -- (msm);
    \draw[arr] (msm) -- (gm);
    \draw[arr] (msm) -- (lm);
    \draw[arr] (lm) -- (mf);
    \draw[arr] (lm) -- (ht);
  \end{tikzpicture}
  \caption{Proposed multiscale analysis architecture.
    Solid boxes are concrete classes; dashed boxes are
    interfaces (concepts or abstract classes).}
  \label{fig:ms-arch}
\end{figure}

% -------------------------------------------------------------------------
\subsection{Core Abstractions}

\paragraph{\texttt{MultiscaleAnalysis}.}
The top-level orchestrator, analogous to \texttt{DynamicAnalysis}
or \texttt{NonlinearAnalysis} in the existing codebase.  It owns
a \texttt{MultiscaleModel} and drives the staggered iteration loop.
Its \texttt{step()} method encapsulates:
\begin{enumerate}
  \item Global macro-scale solve (one time step or load increment).
  \item Kinematics extraction at critical elements.
  \item Parallel local-scale solves (sub-model updates).
  \item Homogenisation and tangent injection.
  \item Coupling convergence evaluation.
\end{enumerate}

\paragraph{\texttt{MultiscaleModel}.}
An aggregation of model types: one \texttt{GlobalModel}
(e.g.\ \texttt{Model<TimoshenkoBeam3D, \ldots>}) and a vector
of \texttt{LocalModelAdapter} instances.  The
\texttt{MultiscaleModel} maintains the transition map
(beam element $\to$ local model index) and provides the interface
for kinematics transfer and result injection.

\paragraph{\texttt{LocalModelAdapter}.}
A type-erased or concept-constrained interface for local-scale
solvers.  Any class satisfying:
\begin{lstlisting}[caption={\texttt{LocalModelAdapter} concept (sketch).}]
template <typename T>
concept LocalModelAdapter = requires(T t,
    const SectionKinematics& kA,
    const SectionKinematics& kB,
    double time)
{
    { t.update_kinematics(kA, kB) };
    { t.solve_step(time) } -> SubModelSolverResult;
    { t.section_forces() } -> Eigen::Vector<double, 6>;
    { t.section_tangent(double) }
        -> Eigen::Matrix<double, 6, 6>;
    { t.revert_state() };
    { t.commit_state() };
};
\end{lstlisting}
can serve as a local model.  This decouples the coupling
framework from the specific solver implementation.

\paragraph{\texttt{ScaleBridgePolicy}.}
A strategy that controls the coupling direction:
\begin{itemize}
  \item \texttt{OneWayCoupling}: downscaling only
    (kinematics $\to$ sub-model, no feedback).
  \item \texttt{TwoWayStaggered}: full staggered iteration
    with tangent injection.
  \item \texttt{TwoWayMonolithic}: assembles a combined
    system (future extension).
\end{itemize}

\paragraph{\texttt{CouplingConvergence}.}
An interface for convergence criteria:
\begin{lstlisting}[caption={Convergence criterion interface.}]
struct CouplingConvergence {
    virtual ~CouplingConvergence() = default;
    virtual bool converged(
        std::span<const Eigen::Matrix<double,6,6>> D_prev,
        std::span<const Eigen::Matrix<double,6,6>> D_curr
    ) const = 0;
};
\end{lstlisting}
Implementations include Frobenius-norm relative tolerance,
energy-norm criteria, and force-equilibrium residuals.

\paragraph{\texttt{RelaxationPolicy}.}
Controls under-relaxation between staggered iterations:
\begin{itemize}
  \item \texttt{ConstantRelaxation}: $\mathbf{D}^{(k+1)} =
    \omega \, \mathbf{D}_{\text{new}} +
    (1 - \omega) \, \mathbf{D}^{(k)}$.
  \item \texttt{AitkenRelaxation}: adaptive $\omega$ from
    successive residual vectors.
  \item \texttt{AndersonAcceleration}: multi-secant acceleration
    using the last $m$ iterates.
\end{itemize}

\paragraph{\texttt{MaterialFactory}.}
Decouples material construction from the solver:
\begin{lstlisting}[caption={Material factory concept.}]
struct ConcreteMaterialFactory {
    virtual Material<ThreeDimensionalMaterial>
        create_concrete() const = 0;
};
struct RebarMaterialFactory {
    virtual Material<UniaxialMaterial>
        create_rebar(double area) const = 0;
};
\end{lstlisting}

\paragraph{\texttt{HomogenisationStrategy}.}
Controls how micro-scale results are upscaled:
\begin{itemize}
  \item \texttt{BoundaryReactionHomogenisation}: current approach
    (sum nodal reactions at Face~B).
  \item \texttt{VolumeAveragingHomogenisation}: classical approach
    (\ref{eq:fe2-homog}).
  \item \texttt{EnergyBasedHomogenisation}: Hill--Mandel energy
    equivalence.
\end{itemize}


% =========================================================================
\section{Code Reusability Assessment}
\label{sec:ms-reuse}

% -------------------------------------------------------------------------
\subsection{Directly Reusable Components}

The following components have clean interfaces and correct
implementations that can be adopted with minimal change:

\begin{enumerate}
  \item \textbf{\texttt{section\_forces\_from\_reactions()}} ---
    Algorithm: assemble $\mathbf{f}_{\mathrm{int}}$ from
    displacement, read Face~B nodal forces, sum into section
    resultants.  Generic and material-independent.

  \item \textbf{Perturbation tangent loop} --- The forward-difference
    scheme in \texttt{compute\_homogenized\_tangent()} is independent
    of the material model; it only requires a callable that maps
    kinematics to section forces.

  \item \textbf{\texttt{SectionKinematics} +
    \texttt{compute\_boundary\_displacements()}} --- The kinematic
    transfer from beam section deformations to 3D displacement BCs
    is well isolated in \texttt{FieldTransfer.hh}.

  \item \textbf{\texttt{MultiscaleSubModel} struct} --- A clean
    data container holding domain, grid, kinematics, and boundary
    conditions.

  \item \textbf{Section override mechanism} ---
    \texttt{BeamElement::set\_homogenized\_tangent/forces()} and
    the \texttt{tangent\_override\_} / \texttt{force\_override\_} in
    \texttt{MaterialSection} implement a clean Strategy/Decorator
    pattern for FE\textsuperscript{2} injection.

  \item \textbf{SNES callbacks} ---
    \texttt{FormResidual()} and \texttt{FormJacobian()} are
    static functions parameterised by a \texttt{Context} struct;
    they can be reused in any PETSc-SNES-based solver.
\end{enumerate}

% -------------------------------------------------------------------------
\subsection{Hard-Coded Dependencies Requiring Refactoring}

\begin{enumerate}
  \item \textbf{Material construction in \texttt{first\_solve()}.}
    \texttt{KoBatheConcrete3D} and \texttt{MenegottoPintoSteel} are
    instantiated directly.  Should be replaced by injected factories.

  \item \textbf{Staggered loop in \texttt{main()}.}
    Approximately 80 lines of coupling algorithm embedded in the
    application driver, including convergence criterion, relaxation
    factor, and ``lagged'' coupling logic.

  \item \textbf{\texttt{extract\_beam\_kinematics()} lambda.}
    Defined locally in \texttt{main\_lshaped\_multiscale.cpp}; should
    be a member of \texttt{MultiscaleModel} or a reusable free
    function.

  \item \textbf{\texttt{SubModelSolver} vs.\
    \texttt{NonlinearSubModelEvolver}.}
    Two classes with $>60\%$ duplicated code (material setup,
    BCs, result extraction).  Should be unified under a common
    interface.

  \item \textbf{\texttt{homogenize()} in
    \texttt{HomogenizedSection.hh}.}
    Assumes uniform stress distribution across the section ---
    not extensible to gradient-based or energy-based
    homogenisation.  Should become a strategy.

  \item \textbf{Geometry coupling to \texttt{PrismaticGrid}.}
    The evolver assumes a prismatic grid with
    \texttt{PrismFace::MinZ/MaxZ} faces.  General RVE geometries
    (e.g.\ from Gmsh imports) would require a different boundary
    identification strategy.
\end{enumerate}


% =========================================================================
\section{Parallelism and HPC Scalability}
\label{sec:ms-hpc}

% -------------------------------------------------------------------------
\subsection{Current Parallelism}

The current implementation exploits parallelism at two levels:
\begin{itemize}
  \item \textbf{BC computation:}
    \texttt{MultiscaleCoordinator::build\_sub\_models()} uses OpenMP
    for the embarrassingly parallel boundary displacement evaluation
    (Phase~2 of the coordinator).

  \item \textbf{Sub-model independence:}  Each sub-model's SNES solver
    uses \texttt{PETSC\_COMM\_SELF} (sequential per sub-model), but
    multiple sub-models can be solved concurrently.
\end{itemize}

% -------------------------------------------------------------------------
\subsection{Scalability Bottlenecks}

\begin{enumerate}
  \item \textbf{Serial perturbation tangent.}
    The six perturbation solves in
    \texttt{compute\_homogenized\_tangent()} are executed
    sequentially.  Since each perturbation is independent (all start
    from the same baseline state), they are trivially parallelisable.

  \item \textbf{Serial sub-model evolution.}
    The sub-model loop in the staggered iteration executes
    sequentially.  For $N_s$ sub-models, this misses a factor-$N_s$
    speedup.

  \item \textbf{Global--local communication.}
    Currently implemented via shared-memory access (kinematics
    extraction and tangent injection occur through direct pointer
    dereference).  Distributed-memory parallelism would require
    explicit communication of kinematics and tangent data.
\end{enumerate}

% -------------------------------------------------------------------------
\subsection{HPC Strategy}

\paragraph{Level 1: Thread-level parallelism.}
Use \texttt{std::jthread} or OpenMP tasks for concurrent sub-model
solves.  Each sub-model owns its own PETSc objects
(\texttt{PETSC\_COMM\_SELF}), so no synchronisation is needed
during the SNES solve.

\paragraph{Level 2: MPI sub-communicators.}
For large-scale analyses with $N_s$ sub-models, partition MPI
ranks into sub-communicators:
\begin{equation}
  \texttt{MPI\_Comm\_split}(\texttt{world}, \;
    \text{rank} \bmod N_s, \; \text{rank}, \; \&\text{sub\_comm}).
\end{equation}
Each sub-communicator drives one sub-model's PETSc solve in
parallel.  The global model runs on a dedicated communicator (or
the full world communicator with the sub-model ranks idle during
the global phase).

\paragraph{Level 3: Perturbation parallelism.}
The six perturbation directions can be distributed across MPI ranks
or OpenMP tasks:
\begin{lstlisting}[caption={Parallelised perturbation tangent.}]
#pragma omp parallel for schedule(static)
for (int j = 0; j < 6; ++j) {
    // Each thread has its own copy of U, imposed_solution
    auto [sp, ok] = perturb_and_solve(j, h_pert);
    if (ok) D_hom.col(j) = (sp - s0) / h_pert;
}
\end{lstlisting}
This requires thread-local copies of the PETSc vectors and SNES
objects, which can be pre-allocated in a pool.

\paragraph{Level 4: GPU offloading.}
For sub-models with $>10{,}000$ degrees of freedom, PETSc's GPU
backend (\texttt{MatSetType(MATAIJCUSPARSE)}) can accelerate the
sparse linear algebra.  This is transparent to the sub-model solver
if the matrix and vector types are set via the options database.

\paragraph{Proposed architecture for HPC:}
\begin{lstlisting}[caption={HPC-aware \texttt{MultiscaleAnalysis}.}]
class MultiscaleAnalysis {
    MPI_Comm global_comm_;         // all ranks
    MPI_Comm macro_comm_;          // global model ranks
    std::vector<MPI_Comm> micro_comms_; // per sub-model
    
    void step(double dt) {
        // 1. Global solve on macro_comm_
        // 2. Broadcast kinematics to micro_comms_
        // 3. Parallel local solves (each on micro_comm_)
        // 4. Gather homogenised responses
        // 5. Inject and check convergence
    }
};
\end{lstlisting}


% =========================================================================
\section{Phased Implementation Plan}
\label{sec:ms-plan}

The refactoring is organised into eight phases, ordered by
dependency and decreasing impact.  Each phase is designed to
maintain a green test suite throughout.

\subsection{Phase 0: Build System Cleanup}
\label{sec:ms-plan-phase0}

\begin{itemize}
  \item \textbf{P0.1:} Remove the 20 duplicate
    \texttt{fall\_n\_seismic\_v2}--\texttt{v21} targets.
    Replace with a single parameterised executable.
    \emph{Impact:} ${\sim}35\%$ build-time reduction.
  \item \textbf{P0.2:} Add precompiled headers (PCH) for Eigen,
    PETSc, VTK, and STL containers.
    \emph{Impact:} ${\sim}15{-}20\%$ reduction.
  \item \textbf{P0.3:} Introduce explicit template instantiations
    for \texttt{Model<>}, \texttt{ContinuumElement<>},
    \texttt{NonlinearAnalysis<>} in dedicated \texttt{.cpp} files,
    with \texttt{extern template} declarations in headers.
    \emph{Impact:} ${\sim}20{-}30\%$ reduction.
  \item \textbf{P0.4:} Modularise \texttt{header\_files.hh} into
    tier-specific headers
    (\texttt{fall\_n/core.hh}, \texttt{fall\_n/materials.hh}, etc.).
\end{itemize}

\subsection{Phase 1: \texttt{LocalModelAdapter} Interface}
\label{sec:ms-plan-phase1}

\begin{itemize}
  \item Define the \texttt{LocalModelAdapter} concept (or abstract
    base class).
  \item Refactor \texttt{NonlinearSubModelEvolver} to satisfy the
    concept.
  \item Deprecate \texttt{SubModelSolver} (stateless variant);
    provide a thin adapter if backward compatibility is required.
  \item All existing tests must pass without modification.
\end{itemize}

\subsection{Phase 2: \texttt{MaterialFactory} Injection}
\label{sec:ms-plan-phase2}

\begin{itemize}
  \item Extract material construction into factory interfaces
    (\texttt{ConcreteMaterialFactory},
    \texttt{RebarMaterialFactory}).
  \item Provide default implementations for KoBatheConcrete3D and
    MenegottoPinto.
  \item Inject the factory into the local model adapter via
    constructor parameter.
  \item Extend with at least one alternative material (e.g.\
    linear-elastic concrete for testing) to validate OCP.
\end{itemize}

\subsection{Phase 3: \texttt{MultiscaleModel} Aggregation}
\label{sec:ms-plan-phase3}

\begin{itemize}
  \item Create a \texttt{MultiscaleModel} class that owns the global
    model and a vector of \texttt{LocalModelAdapter} instances.
  \item Implement the transition map (element ID $\to$ local model
    index).
  \item Move \texttt{extract\_beam\_kinematics()} from the main
    driver into \texttt{MultiscaleModel}.
  \item Provide \texttt{inject\_homogenized\_response()} as a
    member function.
\end{itemize}

\subsection{Phase 4: \texttt{MultiscaleAnalysis} Orchestrator}
\label{sec:ms-plan-phase4}

\begin{itemize}
  \item Create \texttt{MultiscaleAnalysis} class that owns a
    \texttt{MultiscaleModel} and orchestrates the staggered loop.
  \item Move the 80-line staggered iteration from
    \texttt{main\_lshaped\_multiscale.cpp} into
    \texttt{MultiscaleAnalysis::step()}.
  \item Accept \texttt{ScaleBridgePolicy}, \texttt{CouplingConvergence},
    and \texttt{RelaxationPolicy} as constructor parameters.
  \item The main driver reduces to ${\sim}20$ lines of configuration
    plus \texttt{analysis.run()}.
\end{itemize}

\subsection{Phase 5: Coupling Strategy Injection}
\label{sec:ms-plan-phase5}

\begin{itemize}
  \item Implement \texttt{OneWayCoupling} and
    \texttt{TwoWayStaggeredCoupling} as
    \texttt{ScaleBridgePolicy} instances.
  \item Implement \texttt{FrobeniusConvergence},
    \texttt{EnergyConvergence} as
    \texttt{CouplingConvergence} instances.
  \item Implement \texttt{ConstantRelaxation},
    \texttt{AitkenRelaxation} as
    \texttt{RelaxationPolicy} instances.
  \item Verify that switching between strategies requires only
    configuration changes (no recompilation of the coupling core).
\end{itemize}

\subsection{Phase 6: Inter-Sub-Model Parallelism}
\label{sec:ms-plan-phase6}

\begin{itemize}
  \item Parallelise sub-model solves within
    \texttt{MultiscaleAnalysis::step()} using \texttt{std::jthread}
    (shared-memory) or \texttt{MPI\_Comm\_split}
    (distributed-memory).
  \item Parallelise the six perturbation directions in the tangent
    computation.
  \item Benchmark scaling: $N_s = 1, 2, 4, 8$ sub-models on 8
    cores.
\end{itemize}

\subsection{Phase 7: Homogenisation Strategy}
\label{sec:ms-plan-phase7}

\begin{itemize}
  \item Refactor \texttt{homogenize()} into a
    \texttt{HomogenisationStrategy} interface.
  \item Implement boundary-reaction and volume-averaging variants.
  \item Verify consistency: both strategies must yield the same
    section forces for elastic material.
\end{itemize}

\subsection{Phase 8: Header Modularisation and Compilation}
\label{sec:ms-plan-phase8}

\begin{itemize}
  \item Decompose \texttt{header\_files.hh} into module-specific
    aggregation headers.
  \item Create a CMake OBJECT library
    (\texttt{fall\_n\_core}) with explicit template instantiations.
  \item All executables link against the OBJECT library.
  \item Evaluate C++20 modules readiness (GCC~15 experimental
    support).
\end{itemize}


% =========================================================================
\section{Summary}
\label{sec:ms-arch-summary}

The \texttt{fall\_n} multiscale framework successfully implements the
core FE\textsuperscript{2} coupling loop with verified two-way
interaction between beam and continuum scales.  The main contribution
relative to the classical Feyel--Chaboche method is the
\emph{beam--continuum dimensional bridge with boundary-reaction
homogenisation}, which provides a computationally efficient
alternative to full stiffness condensation.

The architectural audit reveals that the current implementation,
while functionally correct, concentrates too many responsibilities
in \texttt{NonlinearSubModelEvolver} and hard-codes material choices
and coupling parameters.  The proposed refactoring into
\texttt{MultiscaleAnalysis} $\to$ \texttt{MultiscaleModel} $\to$
\texttt{LocalModelAdapter} with injected strategies for coupling,
convergence, relaxation, material, and homogenisation will yield a
flexible, extensible, and HPC-ready framework aligned with SOLID
principles and suitable for publication as a reusable library for
the structural engineering community.


% =========================================================================
% References (to be included in the global bibliography)
% =========================================================================
%
% The following references are cited in this chapter.  When a global
% .bib file is used, these entries should be merged there.
%
% @article{Feyel1999,
%   author  = {F. Feyel},
%   title   = {Multiscale {FE2} elastoviscoplastic analysis of composite structures},
%   journal = {Computational Materials Science},
%   volume  = {16},
%   number  = {1--4},
%   pages   = {344--354},
%   year    = {1999},
%   doi     = {10.1016/S0927-0256(99)00077-4}
% }
%
% @article{FeyelandChaboche2000,
%   author  = {F. Feyel and J.-L. Chaboche},
%   title   = {{FE2} multiscale approach for modelling the elastoviscoplastic
%              behaviour of long fibre {SiC/Ti} composite materials},
%   journal = {Computer Methods in Applied Mechanics and Engineering},
%   volume  = {183},
%   number  = {3--4},
%   pages   = {309--330},
%   year    = {2000},
%   doi     = {10.1016/S0045-7825(99)00224-8}
% }
%
% @article{Miehe2002,
%   author  = {C. Miehe and J. Schröder and J. Schotte},
%   title   = {Computational homogenization analysis in finite plasticity.
%              Simulation of texture development in polycrystalline materials},
%   journal = {Computer Methods in Applied Mechanics and Engineering},
%   volume  = {171},
%   number  = {3--4},
%   pages   = {387--418},
%   year    = {1999},
% }
%
% @article{Geers2010,
%   author  = {M. G. D. Geers and V. G. Kouznetsova and W. A. M. Brekelmans},
%   title   = {Multi-scale computational homogenization: Trends and challenges},
%   journal = {Journal of Computational and Applied Mathematics},
%   volume  = {234},
%   number  = {7},
%   pages   = {2175--2182},
%   year    = {2010},
%   doi     = {10.1016/j.cam.2009.08.077}
% }
%
% @article{Kouznetsova2004,
%   author  = {V. G. Kouznetsova and M. G. D. Geers and W. A. M. Brekelmans},
%   title   = {Multi-scale second-order computational homogenization of
%              multi-phase materials: a nested finite element solution
%              strategy},
%   journal = {Computer Methods in Applied Mechanics and Engineering},
%   volume  = {193},
%   number  = {48--51},
%   pages   = {5525--5550},
%   year    = {2004},
%   doi     = {10.1016/j.cma.2003.12.073}
% }
%
% @article{Ghosh2001,
%   author  = {S. Ghosh and K. Lee and S. Moorthy},
%   title   = {Multiple scale analysis of heterogeneous elastic structures
%              using homogenization theory and Voronoi cell finite element
%              method},
%   journal = {International Journal of Solids and Structures},
%   volume  = {32},
%   pages   = {27--62},
%   year    = {1995},
% }
%
% @article{Feyel2003,
%   author  = {F. Feyel},
%   title   = {A multilevel finite element method ({FE2}) to describe the
%              response of highly non-linear structures using generalized
%              continua},
%   journal = {Computer Methods in Applied Mechanics and Engineering},
%   volume  = {192},
%   number  = {28--30},
%   pages   = {3233--3244},
%   year    = {2003},
%   doi     = {10.1016/S0045-7825(03)00348-7}
% }
%
% @article{Schroeder2014,
%   author  = {J. Schröder},
%   title   = {A numerical two-scale homogenization scheme: the {FE2}-method},
%   journal = {Plasticity and Beyond},
%   volume  = {550},
%   pages   = {1--64},
%   year    = {2014},
%   publisher = {Springer},
%   doi     = {10.1007/978-3-7091-1625-8_1}
% }
